\section{The block-diagonal alternating quotient of an imprimitive carrier}
\label{sec:heisenberg-imprimitive-diagonal-carrier}

\paragraph{Author freeze and provenance.}
D-TD42, reserved question Q2ex, is a separate extension of the
commuting-field carrier argument.  Root proposed the block-kernel
dichotomy, transport to every block and the uniform summation below;
\texttt{/root/g34\_eightfolds} supplies its complete proof.  The
terminal-factor analysis has the same exact Guralnick--Saxl,
Schreier and repetition-list inputs as D-TD38.  No source is
assumed for a block stabilizer or terminal coordinate group.
Complete author proof, source-interface, checker and first-output
readings pass.  Independent and root whole-proof reviews also pass.

\paragraph{Actual setup.}
Retain (HCF.1)--(HCF.3): an actual connected finite sphere cover
with coefficient space $V=W\otimes R$, moving-rank sum $2e+6$,
and an actual stable whole nondegenerate projective Heisenberg
factor whose defining action and quotient are
\[
 \begin{gathered}
 J\twoheadrightarrow P\twoheadrightarrow H
            \mathop{\twoheadrightarrow}^{\psi}T\in\{A_n,S_n\},
 \qquad n\ge2^{30},\\
 Q\simeq\mathbf F_p^{2m},\qquad H=P/Q\hookrightarrow\operatorname{Sp}(Q),
 \qquad F=\operatorname{End}_{\mathbf F_pH}(Q)=\mathbf F_q,\\
 q=p^f,\qquad a=\dim_F Q,\qquad af=2m,\qquad
 e=\dim V\ge\dim W=q^{a/2}.
 \end{gathered}
 \tag{HBD.1}
\]
The original $\mathbf F_pH$-module $Q$ is irreducible, so its
$F$-model is absolutely irreducible by D-TD34.  The entire external
partner $R$ is arbitrary.  Put $M=Q\otimes_F\overline{\mathbf F}_p$.

Assume that $M$ is imprimitive, and choose a finest nontrivial
$H$-stable direct-sum block system.  Thus
\[
 M=\bigoplus_{j=1}^{t}U_j,\qquad t\ge2,\qquad
 \dim U_j=b,\qquad a=tb,\qquad
 K=\ker\bigl(H\to\operatorname{Sym}\{U_1,\ldots,U_t\}\bigr).
 \tag{HBD.2}
\]
Finest means that it admits no proper refinement by another
$H$-stable direct-sum system.  Such a choice exists because the
number of nonzero blocks is at most $a$.  Irreducibility makes
$H$ transitive on the blocks, hence makes their dimensions equal.
The blocks are over the algebraic closure; no descent to $F$ or
the prime field is assumed.

\begin{theorem}[D-TD42]
In (HBD.1)--(HBD.2), suppose
\[
                             \psi(K)\supseteq A_n.
 \tag{HBD.3}
\]
Then the actual sphere moving-rank sum $2e+6$ is impossible.
No transitivity of $\ker\psi$ on the blocks is required.
The complementary case $\psi(K)=1$ is not asserted to satisfy
the conclusion of this theorem.
\end{theorem}

\begin{proof}
\paragraph{Step 1: every block retains the actual alternating quotient.}
For each $j$, put
$K_j=\{k\in K:k|_{U_j}=1\}$.  The $K_j$ are normal in $K$,
and their intersection is one, because the original matrix
action is faithful.  The map $\psi|_K$ is onto
$\psi(K)\in\{A_n,S_n\}$.  The coordinate-kernel argument
of D-TD17 gives some $j$ with $K_j\subseteq\ker\psi$.
Here is the full argument.  If every image $\psi(K_j)$ were
nontrivial, each would contain $A_n$.  Inductively the image of
$I_s=\bigcap_{j\le s}K_j$ would contain $A_n$, since
\[
 [I_{s-1},K_s]\le I_s,\qquad
 \psi([I_{s-1},K_s])
        =[\psi(I_{s-1}),\psi(K_s)]\supseteq A_n.
 \tag{HBD.4}
\]
The final trivial intersection contradicts this.  Thus one
coordinate kernel is killed.  For $hU_i=U_j$ one has
$K_j=hK_ih^{-1}$, and hence
$\psi(K_j)=\psi(h)\psi(K_i)\psi(h)^{-1}=1$.
Transitivity of the actual $H$ therefore gives
\[
                         K_j\subseteq\ker\psi
                              \quad\hbox{for every }j.
 \tag{HBD.5}
\]
This uses no transitivity assumption on $\ker\psi$.

Let $H_j$ be the stabilizer of $U_j$.  If $h\in H_j$ acts
as a scalar on $U_j$, then for every $k\in K$ the commutator
$[h,k]$ belongs to $K_j$: it lies in $K$ by normality and acts
trivially on this block.  By (HBD.5), $\psi(h)$ centralizes
$\psi(K)\supseteq A_n$.  The centralizer of the natural
$A_n$ in $A_n$ or $S_n$ is trivial for $n\ge5$, so
$\psi(h)=1$.  Consequently the actual projective block image
has an actual quotient
\[
 P_j=\operatorname{im}(H_j\to\operatorname{PGL}(U_j)),\qquad
 \psi_j:P_j\twoheadrightarrow T_j:=\psi(H_j)\in\{A_n,S_n\}.
 \tag{HBD.6}
\]
It contains $A_n$ because $K\le H_j$.  This argument kills
the full projective block kernel, not merely the elements of
$K$ that vanish on the block.

\paragraph{Step 2: primitive block modules and one fixed terminal coordinate.}
The $H_j$-module $U_j$ is irreducible.  Otherwise a nonzero
proper invariant subspace and its translates over the blocks
would give a proper $H$-submodule of $M$.
It is primitive: an $H_j$-permuted direct-sum decomposition
would transport to a strictly finer $H$-stable system in (HBD.2).
The transport is well defined because different choices of a
map from $U_j$ to another block differ by an element of $H_j$.
In particular $b\ge2$, since (HBD.6) is impossible in degree one.

Fix $U_1$.  Split every factorwise preserved tensor product of
its actual projective image until every factor is terminal.
Strict dimension decrease terminates, leaving no residual factor.
As in the fully supplied construction in D-TD38, irreducibility
and primitivity pass to these factors, and each terminal image
is tensor-indecomposable in the exact Guralnick--Saxl convention.
Its full special-linear inverse image is finite, even in defining
characteristic.  Applying the same coordinate-kernel argument to
(HBD.6) selects one terminal projective image with an actual
surjection to $T_1$.  Write $r\ge2$ for its degree and $G$ for
its finite special-linear lift.  Thus
\[
 U_1=S\otimes Z,\qquad \dim S=r,\qquad \dim Z=b/r,
 \qquad G\twoheadrightarrow T_1.
 \tag{HBD.7}
\]
The decomposition here concerns the block representation over
the algebraic closure.  Neither its factors nor $G$ are required
to descend to $F$ or to occur as subgroups of the original $H$.

Choose $g_j\in H$ with $g_jU_1=U_j$, and transport this entire
choice to each block.  For an operator $h$ fixing $U_j$, the
coordinate analysis is that of $g_j^{-1}hg_j\in H_1$.
Its natural quotient is
$\psi(g_j)^{-1}\psi(h)\psi(g_j)$, so it moves exactly the
same number $k$ of natural letters as $\psi(h)$.
The selected terminal degree and branch are therefore the same
on every fixed block.  They are chosen before the operator $h$,
not selected afresh from its eigenspaces.  If $T_1=A_n$, an
operator with odd natural image cannot fix any block; there is
no missing odd extension in that case.

\paragraph{Step 3: the exhaustive terminal alternatives.}
For a geometric scalar target write $\nu_{\rm alg}$ for
dimension minus largest geometric eigenspace dimension.
If $r\le575$, every quotient-nonidentity coordinate is
nonscalar, giving $\nu_{\rm alg}\ge1$.
If $r\ge576$, apply the exact published Guralnick--Saxl
Theorem7.1 and Remark3 to $G$.\cite{BI:GS}
In its generic branch, every nonscalar coordinate has
$\nu_{\rm alg}>\sqrt r/2$, hence $\nu_{\rm alg}\ge13$.

For clarity, the exceptional-core reduction is algebraic and
requires no sphere cover for $H_1$ or $G$.  Since $r>10$, only
the full natural classical and ordinary deleted cases remain.
The supplied Section7 proof identifies
$S_0=F^*(G/Z(G))$ as simple in these cases.  The stable derived
subgroup $E$ of its central preimage is actual, normal and
quasisimple.  The nonzero $E$-socle is $G$-invariant, so equals
the irreducible module.  Primitivity makes the restriction
homogeneous; the evaluation tensor product is factorwise
preserved, so terminal tensor indecomposability makes its
multiplicity one.  Thus $E$ is absolutely irreducible.

Its centralizer is scalar.  Perfectness kills any automorphism
of $E$ inducing the identity on $E/Z(E)$, since the difference
is a central homomorphism.  Hence
$G/(EZ(G))\hookrightarrow\operatorname{Out}(S_0)$, and the
exact Schreier input makes $G/E$ solvable.
\cite[Theorem2.3, p.~1314]{HD:BabaiLuksSeress1997}
Its actual normal perfect image in the nonsolvable $T_1$ is
therefore $A_n$, forcing $S_0\simeq A_n$.
The complete classical/alternating repetition list in Li--Xia
permits only $A_5,A_6,A_8$ in these named classical families,
so excludes the full natural classical core here.
\cite[Section2.2.1, numbered author p.~13]{HES:LiXia2022}
This is not a claim about arbitrary subgroups of classical groups.

In the remaining deleted case $E$ is the actual ordinary
$A_c$ on its fully deleted module.  A scalar extension has
the same perfect derived group, and $E$ maps injectively onto
$A_n$; hence $c=n$.  After relabeling by a natural automorphism
of $A_n$, each element of $G$ differs by a scalar from the
ordinary deleted operator of its actual image in $T_1$:
they implement the same conjugation on the irreducible core.
Thus $r=n-1$ or $n-2$, and the all-scalar cycle equations of
D-TD35 give, when the natural image moves $k>0$ letters,
\[
             \nu_{\rm alg}\ge\max\{1,k/2-2\}.
 \tag{HBD.8}
\]
Indeed the stronger lower bound is $\operatorname{ind}_n-2$,
and $\operatorname{ind}_n\ge k/2$.  Projective faithfulness
supplies the separate one.  Both estimates hold for every
scalar target in every characteristic, and cover every
compatible odd operator whenever $T_1=S_n$.

\paragraph{Step 4: fixed blocks and moved cycles contribute to the same target.}
Fix an actual $h\in H$ with $\psi(h)\ne1$, and let $k$ be
its moved natural support.  On a fixed block, retain the entire
tensor factor $Z$ of (HBD.7).  The all-characteristic geometric
tensor inequality D-TD27 gives, for every target $\lambda$,
rank at least $(b/r)\delta(k)$, where
\[
 \delta(k)=
 \begin{cases}
 1,&r\le575,\\
 13,&r\ge576\text{ in the generic branch},\\
 \max\{1,k/2-2\},&r=n-1\text{ or }n-2\text{ in the deleted branch}.
 \end{cases}
 \qquad 0<\frac{\delta(k)}r\le\frac12.
 \tag{HBD.9}
\]
The last inequality uses respectively $r\ge2$, $r\ge576$,
and $k\le n$, $r\ge n-2$, $n\ge6$.

For an orbit of $\ell\ge2$ moved blocks, each of dimension
$b$, a $\lambda$-eigenvector is determined by its component
in the first block.  Successive equations determine all other
components and the final equation only imposes more conditions.
Thus the kernel has dimension at most $b$ and the rank at least
\[
            (\ell-1)b\ge\frac{\ell b}{2}
                     \ge\frac{\ell b}{r}\delta(k).
 \tag{HBD.10}
\]
The zero target has zero kernel because $h$ is invertible.
This argument works when $p$ divides $\ell$; it does not
diagonalize a block permutation.

Sum the fixed-block and moved-cycle estimates for that same
target.  Every dimension is counted, giving
\[
 \nu_{\rm alg}(h\mid M)\ge\frac ar\delta(k),\qquad
 \operatorname{rank}_{\mathbf F_p}(1-h)
       =f\operatorname{rank}_{\overline{\mathbf F}_p}(1-h\mid M)
       \ge f\frac ar\delta(k).
 \tag{HBD.11}
\]
There is no hypothesis that $\ker\psi$ is transitive on
blocks, no fixed-block requirement on an individual operator,
and no suppressed multiplicity space.  The field factor $f$
comes only from the original whole $F$-matrix.  A block-field
degree has not been substituted for it.

\paragraph{Step 5: the actual whole-source contradiction.}
D-TD34 gives $q<28900$.  The elementary original-group order
comparison (HPT.3), independent of primitivity, gives $a\ge8192$
at $n\ge2^{30}$.  For $r\le575$, (HBD.11) is greater than14;
in the generic $r\ge576$ branch it is at least13.  The exact
Heisenberg trace ratio is at most $2^{-13/2}<1/90$ on every
nonidentity natural image.

Retain every original inertia period in $J$.  If its quotient
order is $o$, Fourier averaging gives moving rank at least
$e(1-c)(1-1/o)$, where $c=2^{-13/2}$; kernel powers have
trace bound one, and the entire partner cannot increase the
normalized trace.  The actual $T$-quotient cover has genus at
least two and orbifold sum at least $85/42$, as proved in
D-TD02.  Hence
\[
                 e(1-85c)\le252,\qquad e<4536,
                 \qquad e\ge q^{a/2}\ge2^{4096}>4536,
 \tag{HBD.12}
\]
a contradiction.  Only the actual original quotient cover is
used; the stabilizer and terminal groups have no asserted cover.

In the deleted case (HBD.11) gives rank at least one and
at least $k/2-2$.  The small/large support argument of D-TD35
therefore applies to the whole defining action: for
$0<k\le n/16$, the trace ratio is at most
$1/\sqrt2<(15/16)^3\le((n-k)/n)^3$; for $k>n/16$, it is
less than $2^{1-n/64}<n^{-3}$.  Since $a\ge r\ge n-2$,
\[
 e\ge2^{(n-2)/2}>n^3,\qquad
 \frac{|\operatorname{tr}_V(g)|}{e}
 \le\left(\frac{\max\{\operatorname{fix}_n(\psi(g)),1\}}n\right)^3,
 \qquad\psi(g)\ne1\ \Longrightarrow\quad
       \frac{|\operatorname{tr}_V(g)|}{e}<\frac78.
 \tag{HBD.13}
\]
Here $\psi(g)$ is the actual composition from $J$ in (HBD.1).
Reciprocal finite unitary tensor lifts prove these inequalities
on every full original power, with the whole partner retained.

If the original coefficient action is not faithful, its kernel
acts as a scalar on $W$ and hence kills $P,H,T$.  Pass to the
actual connected quotient $C/\ker V$; the faithful coefficient
image has the same moving-rank sum after identity inertia images
are discarded, and the actual $T$-map and all bounds descend.
Accepted HD-FT27 contradicts (HBD.13) and $2e+6$.
This completes every terminal case and proves the theorem.
\end{proof}

\paragraph{Remaining boundary.}
The theorem excludes precisely the branch in which the block-diagonal
kernel of a finest absolute block system has nontrivial alternating
image.  It makes no conclusion here when the entire alternating
quotient factors through the block-permutation image.  The chosen
absolute blocks need not be defined over the original finite field.
No normal alternating subgroup of the original $H$, original
complex character-field equality, polarization comparison or
successful geometric source is inferred.
